\documentclass[12pt,oneside,a4paper,article]{abntex2}
\usepackage[utf8]{inputenc} % Codificação do documento
\usepackage[T1]{fontenc}    % Seleção de código de fonte.
\usepackage[brazil]{babel}  % Idioma do documento
\usepackage{graphicx}       % Inclusão de gráficos
\usepackage{tabularx}       % Tabelas avançadas
\usepackage{amsmath}        % Melhorias em matemática
\usepackage{lipsum}         % Geração de texto dummy
\usepackage{authblk}
\usepackage[num]{abntex2cite} % ou use [num] para citações numéricas
\usepackage[left=3cm,right=2cm,top=3cm,bottom=2cm]{geometry} % Margens
\usepackage{setspace}       % Espaçamento entre linhas
\usepackage{hyperref} % para inserir links

\title{\textbf{Software livre e Ferramentas Open Source}}
\author[1]{Pedro Henrique DOliveira\thanks{pedrohf.oliveira@ucsal.edu.br}}
\author[1]{Guilherme Brito\thanks{⁠guilhermearaujo.brito@ucsal.edu.br}}
\author[1]{Carlito Neto\thanks{carlito.neto@ucsal.edu.br}}
\author[1]{Fernando Chagas\thanks{fernandochagas.santos@ucsal.edu.br}}

\affil[1]{UCSal, Engenharia de Software, Bahia, Brasil}

\date{Março 2026}

\ifthenelse{\equal{\ABNTEXisarticle}{true}}{%
\renewcommand{\maketitlehookb}{}
}{}

 \hypersetup{
colorlinks=false,       % false: boxed links; true: colored links
pdfborder={0 0 0},      % remove as bordas ao redor dos links
}

\renewcommand*{\Authsep}{, }
\renewcommand*{\Authand}{, }
\renewcommand*{\Authands}{, }
\renewcommand*{\Affilfont}{\normalsize\normalfont}
\renewcommand*{\Authfont}{\bfseries}    % make author names boldface    
\setlength{\affilsep}{2em}   % set the space between author and affiliation
\newsavebox\affbox

\begin{document}

\newcommand{\introduction}{
\softwareDefinition
\openSourceSoftwareDefinition
\openSourceToolsDefinition
}

\newcommand{\softwareDefinition}{\par\indent \textbf{Software} é o conjunto de programas, instruções e dados que permitem ao computador executar diferentes tarefas,\cite{david:2001} \textbf{Software Livre e Ferramentas Open Source} são softwares que tem direitos de cópia e distribuição que garantem a sua livre distribuição, uso, modificação e estudo. São representados pela parte lógica do computador, diferentemente do hardware, que corresponde às partes materiais como, processador, placa de video, placa-mãe e circuitos. O software funciona como um conjunto de comandos que orienta o hardware sobre o que deve ser feito, permitindo que o usuário interaja com a máquina. O software é fundamental para o funcionamento dos dispositivos computacionais, pois possibilita que o hardware execute instruções e que os usuários utilizem as máquinas para diversas finalidades.\par \openSourceSoftwareDefinition}

\newcommand{\openSourceSoftwareDefinition}{Diferentemente dos softwares proprietários\cite{fsf:2023}, em que o código-fonte é fechado e controlado por uma empresa ou desenvolvedor, o software livre permite que qualquer pessoa tenha acesso ao código e possa analisá-lo, adaptá-lo ou melhorá-lo conforme suas necessidades. Esse tipo de modelo incentiva a colaboração entre programadores e comunidades ao redor do mundo, contribuindo para o desenvolvimento e aprimoramento contínuo dos programas. \par Além disso, o software livre está presente em diferentes categorias amplamente utilizadas na computação moderna. Entre os principais tipos, destacam-se os sistemas operacionais, como o Linux, que servem como base para a execução de outros programas, os sistemas de gerenciamento de banco de dados, como o MySQL, as ferramentas de desenvolvimento e controle de versão, como o Git,além de aplicações de uso cotidiano como o navegador Mozilla Firefox. \cite{weber:2004}A  ampla adoção dessas ferramentas demonstra como o modelo \textbf{Open Source} se consolidou como um dos pilares da infraestrutura digital global, permitindo inovação contínua e colaboração em larga escala.}

\newcommand{\openSourceToolsDefinition}{ \par \textbf Esse O modelo \textbf{Open Source} de desenvolvimento se baseia em princípios estabelecidos por organizações como a Open Source Initiative e a Free Software Foundation, que definem as liberdades e as estruturas de licenciamento que sustentam os ecossistemas abertos. O desenvolvimento de código aberto se beneficia de contribuições descentralizadas \cite{raymond:2001, fogel:2005}, permitindo que o software evolua de forma mais eficiente em comparação com modelos fechados, destacando-se a importância da colaboração impulsionada pela comunidade para a sustentabilidade desses projetos. Economicamente, o código aberto reduz custos e incentiva a inovação por meio do compartilhamento de conhecimento \cite{lerner:2002}. Exemplos práticos desse paradigma incluem ferramentas amplamente adotadas, como o sistema operacional Linux e Android, sistemas de bancos de dados relacionais como o PostgreSLQ e MySQL que receberam melhorias de performance por colaboradores e pesquisadores externos \cite{repas:2022},demonstram como a colaboração aberta pode resultar em soluções tecnológicas robustas, escaláveis e confiáveis, utilizadas tanto na academia quanto na indústria.}
\newcommand{\developmenttext}{
\developmentIntroduction
\fourPrinciples
\openSourceCorrelation
\globalInfrastructure
\commonTechnologies
}

\newcommand{\developmentIntroduction}{As quatro liberdades do \textbf{Software Livre} ou modelos de desenvolvimento \textbf{Open Source} são princípios que garantem aos usuários controle e autonomia sobre os programas que utilizam. Essas liberdades permitem que qualquer pessoa utilize, estude, modifique, e compartilhe o software sem restrições impostas pelo desenvolvedor original. Elas são consideradas a base do conceito de software livre e incentivam a colaboração, o aprendizado e o desenvolvimento tecnológico. \par \fourPrinciples}

\newcommand{\fourPrinciples}{
A \textbf{Usar o Software}. Isso significa que o usuário pode utilizar o software e executa-lo da forma que desejar.Isso significa que o usuário pode utilizar o software da forma que desejar, seja para fins pessoais, acadêmicos, profissionais ou comerciais, sem precisar pedir autorização ao criador do programa. Essa liberdade garante que o software possa ser usado em diferentes contextos, sem limitações impostas pelo desenvolvedor.

A \textbf{Estudar o Software}, verificar como o programa funciona e adaptá-lo às próprias necessidades. Para que isso seja possível, é necessário que o \textbf{código-fonte} do software esteja disponível para os usuários. Dessa forma, programadores e estudantes podem analisar como o programa foi desenvolvido, entender sua estrutura e modificar partes do código para melhorar seu funcionamento ou adicionar novas funcionalidades.

A \textbf{Distribuir o Software}, redistribuir cópias do software permite que os usuários compartilhem o programa com outras pessoas, seja gratuitamente ou até mesmo cobrando por essa distribuição. Essa liberdade facilita o acesso à tecnologia e contribui para que mais pessoas possam utilizar e conhecer o software.

A \textbf{Modificar o Software},e disponibilizar essas modificações para outros usuários. Isso significa que, além de alterar o programa para uso próprio, o usuário também pode compartilhar versões modificadas com a comunidade. Dessa forma, melhorias feitas por um usuário podem beneficiar muitos outros, promovendo a evolução constante do software.}

\newcommand{\openSourceCorrelation}{ Atualmente, \textbf{Software Livre e Ferramentas Open Source} ocupam um papel central em diversas camadas da tecnologia da informação, estando presente desde sistemas operacionais até ferramentas de desenvolvimento e gerenciamento de projetos. Esse crescimento está diretamente relacionado aos seus princípios fundamentais, definidos pela Free Software Foundation e pela Open Source Initiative, que garantem a liberdade de uso, modificação e redistribuição do código, permitindo sua adaptação a diferentes contextos e necessidades \cite{fsf:2023, osi:2023}. Esse modelo colaborativo favorece a inovação contínua, contrastando com o modelo fechado de desenvolvimento tradicional \cite{raymond:2001}. Na prática, tecnologias amplamente utilizadas como Linux, MySQL e GitHub evidenciam como o software livre sustenta grande parte da infraestrutura digital moderna.

No contexto de ferramentas de gerenciamento de projetos, estudos demonstram que soluções open source, como OpenProject, ProjectLibre e Redmine, oferecem funcionalidades robustas como controle de tarefas, cronogramas, gestão de recursos e acompanhamento de progresso sendo comparáveis às ferramentas proprietárias como JIRA e Microsoft Project \cite{abramova:2016}. Ferramentas proprietárias tem uma tendência a apresentar interfaces mais intuitivas e melhor experiência de usuário, as soluções open source conseguem atender aos requisitos essenciais de gerenciamento de projetos, muitas vezes com maior flexibilidade e possibilidade de customização. Além disso o modelo open source também apresenta vantagens econômicas, reduzindo custos e incentivando a participação de comunidades no desenvolvimento e evolução dos sistemas\cite{lerner:2002}.

Por fim, observa-se que a adoção de software livre não se limita ao uso individual, mas se estende à infraestrutura global, sendo amplamente empregado em servidores, redes e dispositivos móveis, como no caso do sistema Android. Dessa forma, mesmo que usuários finais frequentemente optem por soluções proprietárias pela facilidade de uso, o software livre permanece como base tecnológica essencial, promovendo independência, transparência e inovação contínua no ecossistema digital.

}

\newcommand{\globalInfrastructure}{
O software de código aberto tornou-se um componente fundamental da infraestrutura digital moderna, sustentando a vasta maioria dos sistemas tecnológicos utilizados na indústria e nos serviços públicos. Estudos recentes realizados na Europa em Organizações de Setor Publico, indicam que mais de 90\% da infraestrutura digital global depende de componentes de código aberto, destacando seu papel crucial na viabilização de sistemas escaláveis, interoperáveis e colaborativos \cite{linaker:2026}. Além disso, o código aberto atua como o “tecido conjuntivo” da economia digital global, dando suporte a serviços essenciais como bancos, saúde e computação em nuvem \cite{forsgren:2021}. Essa ampla adoção reforça sua importância como base estratégica para ecossistemas tecnológicos resilientes e seguros \commonTechnologies.
}

\newcommand{\commonTechnologies}{ \par O papel central do software de código aberto na infraestrutura global se reflete diretamente na ampla adoção de tecnologias consagradas utilizadas no desenvolvimento de software moderno. Plataformas como o Android, baseado no kernel Linux, demonstram como componentes de código aberto impulsionam bilhões de dispositivos em todo o mundo, particularmente no ecossistema móvel. Da mesma forma, ferramentas como o Git\cite{git:2026} se tornaram o padrão para controle de versão distribuído, possibilitando o desenvolvimento colaborativo em escala entre organizações e comunidades abertas. Tecnologias de conteinerização como o Docker\cite{docker:2026} ilustram ainda mais esse impacto, fornecendo ambientes leves e portáteis que suportam aplicações escaláveis e nativas da nuvem. \par Mesmo em ambientes predominantemente proprietários, como o Microsoft Windows, há uma crescente integração com tecnologias de código aberto, incluindo suporte nativo para ferramentas e fluxos de trabalho de desenvolvimento baseados em Linux (WSL). Esses exemplos reforçam as conclusões que o software de código aberto não é apenas fundamental em um nível infraestrutural abstrato, mas também está concretamente incorporado às ferramentas, plataformas e ecossistemas cotidianos que sustentam o mundo moderno}
\newcommand{\conclusion}{ Pode-se observar que o \textbf{Software Livre e Ferramentas Open Source} são pilares para o desenvolvimento de software da atual, contribuindo para a infraestrutura global, acessibilidade e disponibilidade de acesso a ferramentas que ajudam ou melhoram as atividades. A facil colaboração e uso gerada pelos principios e diretrizes deste modelo de desenvolvimento, gera um impacto massivo na econômia e geração de tecnologia fazendo com que a base de grandes empresas e estruturas como a internet sejam alimentadas por estas ferramentas. }
\newcommand{\resume}{Este trabalho apresenta uma visão geral sobre o conceito de \textbf{Software Livre e Ferramentas Open Source} abordando a importância no contexto tecnológico atual e modelo de desenvolvimento de software. São discutidos os princípios fundamentais do software livre, a liberdade de uso, estudo, modificação e distribuição, destacando seu impacto na democratização do acesso à tecnologia e no incentivo à inovação colaborativa. São discutidos exemplos, evidenciando como este tipo de solução contribui para o aumento da produtividade, redução de custos e fortalecimento de comunidades de desenvolvimento. Por fim, ressalta-se o papel estratégico destes modelos, promovendo flexibilidade, transparência e independência tecnológica.}
\newcommand{\questions}{

\noindent\textbf{O que são Software Livre e Ferramentas Open Source?} \\
São modelos de software que podem ser modificados, usados, estudados e distribuídos de forma livre.

\vspace{0.5em}

\noindent\textbf{Quais exemplos de Software Livre e Ferramentas Open Source?} \\
Linux, Git.

\vspace{0.5em}

\noindent\textbf{Quais as 4 diretrizes do Software Livre?} \\
- Usar o software \\
- Modificar o software \\
- Estudar o software \\
- Distribuir o software

\vspace{0.5em}

\noindent\textbf{Por que utilizar Software Livre e Ferramentas Open Source?} \\
Facilidade no desenvolvimento e baixo custo.

\vspace{0.5em}

\noindent\textbf{Qual a importância do Software Livre e Ferramentas Open Source?} \\
Estão presentes na maioria das tecnologias e sistemas utilizados globalmente, sendo parte fundamental da infraestrutura digital.

}

\begin{center}
    \includegraphics[width=0.3\textwidth]{ucsal_logo.png} 
\end{center}

% {\let\newpage\relax\maketitle}

\begin{resumoumacoluna}
\resume
\par\vspace{\baselineskip}
\raggedright
\textbf{Palavras-chaves}: software open source, software open source definitions, open source tools, open source tools definition.
\end{resumoumacoluna}

\textual

\section{Introdução}
\subsection{Software Livre e Ferramentas Open Source}
\softwareDefinition
\subsection{Importância}
\openSourceToolsDefinition


\section{Desenvolvimento}
\subsection{Principios do Software Livre}
\developmentIntroduction
\subsection{Contexto Atual}
\openSourceCorrelation
\subsection{Infraestrutura Global}
\globalInfrastructure

\section{Conclusão}
\conclusion

\section{Perguntas Sobre o Tema}
\questions

\bibliography{references} 
\end{document}
